\documentclass[10pt,letterpaper]{article}
\usepackage[utf8]{inputenc}\usepackage[T1]{fontenc}\usepackage{lmodern}
\usepackage[spanish,es-noshorthands]{babel}
\usepackage[top=0.8cm,bottom=0.8cm,left=1.0cm,right=1.0cm,nohead,nofoot]{geometry}
\usepackage{multicol,amsmath,graphicx,xcolor,array,enumitem,tabularx,booktabs}
\usepackage{pgfplots}\pgfplotsset{compat=1.17}\usetikzlibrary{patterns}
\definecolor{dorado}{HTML}{B8975A}
\definecolor{guinda}{HTML}{6F1D46}\definecolor{tinta}{HTML}{000000}
\pagestyle{empty}\setlength{\parindent}{0pt}\setlength{\columnsep}{0.55cm}
\setlength{\columnseprule}{0.3pt}
\newlength{\wbola}\settowidth{\wbola}{\scriptsize$\bigcirc$}
\begin{document}\fontsize{8.1}{9.15}\selectfont\color{tinta}\columnseprulecolor{\color{tinta}}
\begingroup\color{tinta}
\begin{minipage}{0.09\textwidth}\includegraphics[height=1.15cm]{ipn-escudo.png}\end{minipage}%
\begin{minipage}{0.70\textwidth}\centering
{\color{guinda}\textbf{\large Instituto Politécnico Nacional · ESCA Unidad Santo Tomás}}\\[-1pt]
{\small Laboratorio: Gestión de datos, comportamiento y perfiles del consumidor digital}\\[1pt]
{\normalsize\textbf{Examen de la Unidad I (Capítulos 1 a 3)}} \hfill {\small Profesor: \textbf{Dr. Aboud Barsekh-Onji}}
\end{minipage}%
\begin{minipage}{0.19\textwidth}\centering
{\color{guinda}\fbox{\parbox[c][1.0cm][c]{1.3cm}{\centering\textbf{\scriptsize VERSIÓN}\\[-2pt]{\Huge\textbf{B}}}}}
\end{minipage}
\vspace{2pt}\hrule\vspace{3pt}
{\small\textbf{Nombre:} \rule{7.2cm}{0.3pt} \hfill \textbf{Grupo:} \rule{1.6cm}{0.3pt} \hfill \textbf{Fecha:} \rule{2.2cm}{0.3pt} \hfill \textbf{Calif.:} \rule{1.5cm}{0.3pt}}\\[3pt]
{\footnotesize\textbf{Duración: 50 min.} 30 reactivos de opción múltiple con \textbf{una sola} respuesta correcta. Contesta en la \textbf{hoja de respuestas única} (es la misma para las tres versiones): rellena por completo un solo círculo por reactivo y \textbf{anota la versión que te tocó} (la letra de la esquina de este examen). Cada acierto vale $10/30$ de punto. Algunos reactivos piden analizar las figuras y resultados de arriba. Lee con atención cada enunciado.}
\vspace{2pt}\hrule\vspace{4pt}
\endgroup
\begingroup\tiny\pgfplotsset{every axis/.append style={tick label style={font=\tiny},label style={font=\tiny},title style={font=\tiny},
  every axis plot/.append style={line width=0.6pt}, xlabel style={yshift=2pt}, ylabel style={yshift=-2pt}}}
\noindent\begin{minipage}[t]{0.215\textwidth}\centering
\begin{tikzpicture}[baseline=(current bounding box.north)]\begin{axis}[width=4.5cm,height=4.1cm,xlabel={periodo},ylabel={tasa de adopción},xmin=0,xmax=20,ymin=0,ymax=0.16,
 axis lines*=left,ytick={0,0.05,0.1,0.15},xtick={0,5,10,15,20}]
\addplot[domain=0:20,samples=60,black,dashed]{(0.18)^2*exp(-0.18*x)/0.08/(1+0.10/0.08*exp(-0.18*x))^2};
\addplot[domain=0:20,samples=60,black]{(0.41)^2*exp(-0.41*x)/0.03/(1+0.38/0.03*exp(-0.41*x))^2};
\addplot[domain=0:20,samples=60,black,dotted,line width=1pt]{(0.56)^2*exp(-0.56*x)/0.01/(1+0.55/0.01*exp(-0.56*x))^2};
\node[font=\tiny] at (axis cs:2.6,0.097) {A}; \node[font=\tiny] at (axis cs:4.5,0.125) {B}; \node[font=\tiny] at (axis cs:11.2,0.14) {C};
\end{axis}\end{tikzpicture}\\[-2pt]{\tiny\textbf{Figura 1.} Tasa de adopción (Bass).\\ A: $p{=}.08,q{=}.10$; B: $p{=}.03,q{=}.38$; C: $p{=}.01,q{=}.55$}
\end{minipage}\hfill
\begin{minipage}[t]{0.285\textwidth}\centering
\vspace*{0.5cm}{\tiny\renewcommand{\arraystretch}{1.1}\setlength{\tabcolsep}{2.5pt}\begin{tabular}{@{}lcccc@{}}\toprule
Caso & Arquet. & Traslape & Silueta & \% ambig.\\\midrule
A & 3 & base & 0.58 & 14.6\\ B & 3 & mayor & 0.41 & 36.7\\ C & 4 & base & 0.58 & 21.6\\\bottomrule\end{tabular}}\\[1pt]
\begin{tikzpicture}[baseline=(current bounding box.north)]\begin{axis}[xbar stacked,width=5.0cm,height=2.5cm,bar width=5pt,xmin=0,xmax=100,symbolic y coords={Cl.4,Cl.3,Cl.2,Cl.1},ytick=data,
 xlabel={pertenencia (\%)},axis lines*=left,enlarge y limits=0.22]
\addplot[fill=black!25,draw=black!60] coordinates{(92,Cl.1)(78,Cl.2)(48,Cl.3)(15,Cl.4)};
\addplot[fill=black!55,draw=black!70] coordinates{(5,Cl.1)(12,Cl.2)(44,Cl.3)(20,Cl.4)};
\addplot[fill=white,draw=black!70,pattern=north east lines] coordinates{(3,Cl.1)(10,Cl.2)(8,Cl.3)(65,Cl.4)};
\end{axis}\end{tikzpicture}\\[-2pt]{\tiny\textbf{Figura 2.} Corridas del laboratorio (arriba) y pertenencias de cuatro clientes (abajo, hipotéticas): segmento 1 gris claro, 2 gris oscuro, 3 rayado}
\end{minipage}\hfill
\begin{minipage}[t]{0.215\textwidth}\centering
\begin{tikzpicture}[baseline=(current bounding box.north)]\begin{axis}[ybar,width=4.5cm,height=4.1cm,bar width=8pt,symbolic x coords={18--24,25--34,35--44,55--64,65+},xtick=data,ymin=0,ymax=7,
 xlabel={edad (años)},ylabel={horas diarias de internet},axis lines*=left,enlarge x limits=0.15,nodes near coords,nodes near coords style={font=\tiny},ytick={0,2,4,6},x tick label style={font=\tiny,rotate=45,anchor=east}]
\addplot[fill=black!40,draw=black!70] coordinates{(18--24,5.7)(25--34,5.6)(35--44,4.7)(55--64,3.2)(65+,3.0)};
\end{axis}\end{tikzpicture}\\[-2pt]{\tiny\textbf{Figura 3.} Uso diario de internet por edad en México (ENDUTIH 2024; no se muestra 45--54)}
\end{minipage}\hfill
\begin{minipage}[t]{0.255\textwidth}\centering
\begin{tikzpicture}[baseline=(current bounding box.north)]\begin{axis}[ybar,width=4.9cm,height=4.1cm,bar width=11pt,symbolic x coords={Mujeres,Hombres,Total},xtick=data,ymin=0,ymax=100,
 ylabel={\% con producto financiero formal},axis lines*=left,enlarge x limits=0.3,nodes near coords,nodes near coords style={font=\tiny},ytick={0,25,50,75,100}]
\addplot[fill=black!40,draw=black!70] coordinates{(Mujeres,72.8)(Hombres,80.9)(Total,76.5)};
\end{axis}\end{tikzpicture}\\[-2pt]{\tiny\textbf{Figura 4.} Personas con un producto financiero formal (ENIF 2024, INEGI/CNBV)}
\end{minipage}
\endgroup
\vspace{2pt}\hrule\vspace{3pt}
\begin{multicols}{2}
\noindent\textbf{1.} En 2025 el comercio electrónico en México creció 19.2\,\% anual, mientras que el uso de internet pasó de 83.1\,\% a 86.1\,\% de la población. ¿Qué lectura sugiere esa comparación?
\vspace{-3pt}
\begin{enumerate}[label=\alph*),leftmargin=1.1em,itemsep=0pt,topsep=1pt,parsep=0pt,partopsep=0pt]
\item Buena parte del crecimiento viene de personas ya conectadas que compran más, no solo de personas nuevas que se conectan.
\item El crecimiento del comercio electrónico se explica casi por completo por nuevos usuarios que se conectaron a internet durante ese año.
\item Ambas cifras miden en realidad el mismo fenómeno, por lo que sus tasas de crecimiento tendrían que ser prácticamente iguales.
\end{enumerate}
\vspace{2pt}
\noindent\textbf{2.} Según la Figura 1, ¿en qué periodo ocurre aproximadamente el pico de adopción del canal B ($p=0.03$, $q=0.38$)?
\vspace{-3pt}
\begin{enumerate}[label=\alph*),leftmargin=1.1em,itemsep=0pt,topsep=1pt,parsep=0pt,partopsep=0pt]
\item Alrededor del periodo 1, porque la curva parte de su valor máximo y luego solo disminuye.
\item Alrededor del periodo 6, que coincide con $t^\ast=\ln(q/p)/(p+q)\approx 6.2$ que calcula el script.
\item Alrededor del periodo 12, cuando la adopción acumulada ya casi alcanza todo el mercado potencial.
\end{enumerate}
\vspace{2pt}
\noindent\textbf{3.} Una consultora afirma: «el consumidor 4.0 ya es mayoritario en México porque así lo establece \textit{Marketing 4.0}». ¿Cuál es la crítica más sólida a ese argumento?
\vspace{-3pt}
\begin{enumerate}[label=\alph*),leftmargin=1.1em,itemsep=0pt,topsep=1pt,parsep=0pt,partopsep=0pt]
\item La afirmación es inválida, porque \textit{Marketing 5.0} y \textit{6.0} ya reemplazaron por completo a la versión 4.0 en todos los mercados.
\item La serie \textit{Marketing X.0} es una periodización editorial, no un modelo validado: la afirmación requiere evidencia empírica.
\item La afirmación es válida, porque la serie de Kotler es la teoría académica con mayor revisión por pares sobre la evolución del consumidor digital.
\end{enumerate}
\vspace{2pt}
\noindent\textbf{4.} El capítulo señala que, según fuentes de análisis jurídico, la nueva LFPDPPP admitiría el consentimiento tácito como regla general y a la vez exigiría que sea libre, específico e informado, y no resuelve esa tensión. ¿Por qué?
\vspace{-3pt}
\begin{enumerate}[label=\alph*),leftmargin=1.1em,itemsep=0pt,topsep=1pt,parsep=0pt,partopsep=0pt]
\item Porque el libro considera que el consentimiento no es relevante para la gestión de datos del consumidor digital.
\item Porque la ley es contradictoria y, por esa razón, carece de validez para regular el tratamiento de datos personales.
\item Porque viene de fuentes secundarias no confirmadas contra el texto legal; una postura firme sin verificarlo carecería de sustento.
\end{enumerate}
\vspace{2pt}
\noindent\textbf{5.} Un blog afirma que «existen siete tipos de consumidor digital» sin explicar autoría ni método. ¿Cómo debe tratarlo un analista, según el capítulo?
\vspace{-3pt}
\begin{enumerate}[label=\alph*),leftmargin=1.1em,itemsep=0pt,topsep=1pt,parsep=0pt,partopsep=0pt]
\item Como una tipología sin respaldo académico; lo que sí lo tiene es el método de \textit{clustering} aplicado a datos propios.
\item Como equivalente a \textit{El Cubo NORISO}, que es una tipología de perfiles validada mediante revisión por pares.
\item Como una tipología científica, siempre que el blog tenga muchos seguidores y sea citado con frecuencia en el medio.
\end{enumerate}
\vspace{2pt}
\noindent\textbf{6.} Respecto de k-means y \textit{fuzzy c-means}, ¿qué relación es correcta?
\vspace{-3pt}
\begin{enumerate}[label=\alph*),leftmargin=1.1em,itemsep=0pt,topsep=1pt,parsep=0pt,partopsep=0pt]
\item Son equivalentes cuando el exponente difuso vale 2, por lo que siempre producen exactamente las mismas asignaciones.
\item La difusa es un caso particular de k-means en el que se permiten pertenencias parciales a varios segmentos a la vez.
\item k-means es el caso límite en que toda la pertenencia va a un segmento ($u_{ik}\in\{0,1\}$); la difusa es la general.
\end{enumerate}
\vspace{2pt}
\noindent\textbf{7.} La Figura 4 muestra el porcentaje de mujeres y de hombres con un producto financiero formal (ENIF 2024). Si un banco define «cliente de alto valor» a partir de los productos que sus clientes tienen contratados, ¿qué riesgo advierte el capítulo?
\vspace{-3pt}
\begin{enumerate}[label=\alph*),leftmargin=1.1em,itemsep=0pt,topsep=1pt,parsep=0pt,partopsep=0pt]
\item Que la brecha demuestre que las mujeres no usan servicios financieros y que por ello deban excluirse del análisis.
\item Ninguno: como la definición usa solo números, sus resultados son neutrales respecto del género y de cualquier otra característica.
\item Que la definición herede el sesgo social de los datos de origen y subrepresente a las mujeres sin que ningún paso técnico lo haga explícito.
\end{enumerate}
\vspace{2pt}
\noindent\textbf{8.} La Figura 2 (parte inferior) muestra la pertenencia difusa de cuatro clientes a tres segmentos (datos hipotéticos). ¿Qué cliente ilustra mejor el problema de la asignación única y por qué?
\vspace{-3pt}
\begin{enumerate}[label=\alph*),leftmargin=1.1em,itemsep=0pt,topsep=1pt,parsep=0pt,partopsep=0pt]
\item El cliente 1: su pertenencia de 92\,\% a un segmento muestra que k-means siempre fuerza una casilla que no corresponde.
\item El cliente 3: su mayor pertenencia (48\,\%) no llega a 60\,\%; k-means lo asigna a un segmento sin avisar que es una mezcla.
\item El cliente 4: tiene 35\,\% repartido entre los otros dos segmentos, por lo que es el más ambiguo de los cuatro.
\end{enumerate}
\vspace{2pt}
\noindent\textbf{9.} En una plataforma de anuncios, el equipo define audiencias personalizadas, públicos similares (\textit{lookalike}) y exclusiones. Esto muestra que, en el entorno digital, el mercado meta:
\vspace{-3pt}
\begin{enumerate}[label=\alph*),leftmargin=1.1em,itemsep=0pt,topsep=1pt,parsep=0pt,partopsep=0pt]
\item Solo se declara en el plan de mercadotecnia; la plataforma no interviene en su definición operativa ni en su ejecución.
\item Deja de ser una decisión estratégica de la organización, porque el algoritmo de la plataforma lo elige por sí solo.
\item Además de declararse en el plan, se configura como parámetro, y un error de configuración desperdicia presupuesto en horas.
\end{enumerate}
\vspace{2pt}
\noindent\textbf{10.} Una persona es «ocasional» para su plataforma de streaming y «de alto compromiso» para su banco. ¿Cómo se interpreta?
\vspace{-3pt}
\begin{enumerate}[label=\alph*),leftmargin=1.1em,itemsep=0pt,topsep=1pt,parsep=0pt,partopsep=0pt]
\item El banco usa un algoritmo más preciso, por lo que su etiqueta debe prevalecer también en la plataforma de streaming.
\item No hay contradicción: cada negocio mide un comportamiento distinto y el perfil depende del negocio, no es un atributo fijo.
\item Una de las dos plataformas se equivocó: el perfil verdadero de la persona es único y portable entre todos los negocios.
\end{enumerate}
\vspace{2pt}
\noindent\textbf{11.} Una cadena de farmacias detecta que el 60\,\% de los abandonos de carrito ocurre al pedir un segundo teléfono, sin que ningún cliente se haya quejado. Este hallazgo proviene principalmente de una señal:
\vspace{-3pt}
\begin{enumerate}[label=\alph*),leftmargin=1.1em,itemsep=0pt,topsep=1pt,parsep=0pt,partopsep=0pt]
\item De contexto: la determinan la hora, el canal y el dispositivo en que ocurre la interacción, no el paso del proceso.
\item Conductual: lo que el consumidor hace (dónde abandona), observable aunque nadie lo declare de forma explícita.
\item Declarada: los clientes expresaron su inconformidad de forma explícita mediante reseñas, encuestas o comentarios en redes.
\end{enumerate}
\vspace{2pt}
\noindent\textbf{12.} En la tabla de la Figura 2, del Caso A al Caso C (de 3 a 4 arquetipos) la silueta casi no cambia (0.58), pero el porcentaje de clientes ambiguos pasa de 14.6\,\% a 21.6\,\%. ¿Qué enseña esa comparación?
\vspace{-3pt}
\begin{enumerate}[label=\alph*),leftmargin=1.1em,itemsep=0pt,topsep=1pt,parsep=0pt,partopsep=0pt]
\item Que la silueta de la segmentación rígida por sí sola no revela cuántos clientes quedaron en frontera; la pertenencia difusa sí lo muestra.
\item Que la segmentación difusa falló al aumentar los arquetipos, porque ahora marca como ambiguos a más clientes.
\item Que agregar un cuarto arquetipo no afecta a los clientes, porque la silueta se mantuvo sin cambios.
\end{enumerate}
\vspace{2pt}
\noindent\textbf{13.} Un equipo decide analizar solo tableros de métricas y descarta las entrevistas con clientes por considerarlas «subjetivas». Desde la doble naturaleza del consumidor digital, su error es:
\vspace{-3pt}
\begin{enumerate}[label=\alph*),leftmargin=1.1em,itemsep=0pt,topsep=1pt,parsep=0pt,partopsep=0pt]
\item Tratar al consumidor solo como persona vivida y quedar sin capacidad de operar a la escala de un mercado digital.
\item Ninguno: los datos conductuales son siempre más objetivos que cualquier fuente de contexto humano.
\item Tratar al consumidor solo como dato estructurado y perder el contexto humano de la decisión.
\end{enumerate}
\vspace{2pt}
\noindent\textbf{14.} Un equipo asegura que su decisión es correcta porque sus datos están limpios y se recolectaron con base legal. Según el capítulo, esa condición es:
\vspace{-3pt}
\begin{enumerate}[label=\alph*),leftmargin=1.1em,itemsep=0pt,topsep=1pt,parsep=0pt,partopsep=0pt]
\item Suficiente: un dato limpio y legal garantiza por sí solo una buena decisión de mercadotecnia.
\item Innecesaria: lo decisivo es la intuición del equipo, y los datos solo ilustran lo que ya se había decidido.
\item Necesaria pero no suficiente: además hace falta un método que traduzca el dato en decisión sin engañarse a sí mismo.
\end{enumerate}
\vspace{2pt}
\noindent\textbf{15.} Una app bancaria tiene una interfaz rápida y fácil de usar, pero el servicio al cliente detrás de ella es lento e inconsistente. Esto describe:
\vspace{-3pt}
\begin{enumerate}[label=\alph*),leftmargin=1.1em,itemsep=0pt,topsep=1pt,parsep=0pt,partopsep=0pt]
\item Una CX sólida, porque la facilidad de uso de la app define por sí sola la experiencia completa del cliente.
\item Una UX sólida dentro de una CX deficiente: la interfaz es solo una parte de la experiencia total del cliente.
\item Una UX deficiente, porque la lentitud del servicio es un problema de diseño de la interfaz.
\end{enumerate}
\vspace{2pt}
\noindent\textbf{16.} Un despacho afirma que el INAI sigue siendo la autoridad de protección de datos personales de los particulares en México. Según la cronología del capítulo, ¿qué se puede responder?
\vspace{-3pt}
\begin{enumerate}[label=\alph*),leftmargin=1.1em,itemsep=0pt,topsep=1pt,parsep=0pt,partopsep=0pt]
\item Que el INAI fue extinguido por la reforma de simplificación orgánica y la Secretaría Anticorrupción y Buen Gobierno asumió la protección de datos de particulares.
\item Que ninguna autoridad asumió esas funciones, por lo que hoy la ley no tiene quien vigile su cumplimiento.
\item Que es correcto: la nueva LFPDPPP de 2025 reforzó al INAI y le dio más atribuciones que antes.
\end{enumerate}
\vspace{2pt}
\noindent\textbf{17.} La Figura 3 muestra las horas diarias de uso de internet por grupo de edad en México (ENDUTIH 2024): caen de forma casi monótona al aumentar la edad. ¿Qué advertencia metodológica se desprende para quien segmenta?
\vspace{-3pt}
\begin{enumerate}[label=\alph*),leftmargin=1.1em,itemsep=0pt,topsep=1pt,parsep=0pt,partopsep=0pt]
\item Saltar a variables psicográficas sin controlar edad y zona construye categorías elegantes sobre una base demográfica ya explicativa.
\item Que las variables demográficas son irrelevantes y conviene segmentar únicamente con variables psicográficas de estilo de vida.
\item Que la edad explica por completo el uso de internet, por lo que ya no es necesario segmentar a los consumidores.
\end{enumerate}
\vspace{2pt}
\noindent\textbf{18.} Una empresa lanza un solo producto y le dedica una gran campaña para convencer a todo el mercado de que lo compre. Con la distinción de Wendell Smith (1956), esa estrategia es:
\vspace{-3pt}
\begin{enumerate}[label=\alph*),leftmargin=1.1em,itemsep=0pt,topsep=1pt,parsep=0pt,partopsep=0pt]
\item Mercado meta digital, porque la empresa eligió un único producto y una única campaña para dirigirse a su audiencia.
\item Segmentación de mercado, porque la empresa dirige su campaña a un mercado heterogéneo y busca cubrir todas las preferencias.
\item Diferenciación de producto: adapta el mercado al producto; la segmentación haría lo contrario y adaptaría el producto a subgrupos.
\end{enumerate}
\vspace{2pt}
\noindent\textbf{19.} En una encuesta, 70\,\% de los consumidores dice preferir productos sostenibles, pero esos productos son solo 8\,\% de las ventas (datos hipotéticos). ¿Qué explicación es coherente con el capítulo?
\vspace{-3pt}
\begin{enumerate}[label=\alph*),leftmargin=1.1em,itemsep=0pt,topsep=1pt,parsep=0pt,partopsep=0pt]
\item Los consumidores mienten deliberadamente en las encuestas, por lo que ningún dato declarado tiene valor analítico para la mercadotecnia.
\item Brecha actitud-conducta: la actitud declarada predice mal la compra real, donde pesan el precio, la conveniencia y la confianza en la etiqueta.
\item La sostenibilidad no influye en la decisión de compra, por lo que conviene eliminarla de los análisis de mercado de la empresa.
\end{enumerate}
\vspace{2pt}
\noindent\textbf{20.} ¿Qué cambió realmente con la segmentación digital respecto de la tradicional?
\vspace{-3pt}
\begin{enumerate}[label=\alph*),leftmargin=1.1em,itemsep=0pt,topsep=1pt,parsep=0pt,partopsep=0pt]
\item Encareció la variable conductual, porque ahora exige procesar más datos de los que aportaba una encuesta tradicional.
\item No creó variables nuevas: abarató la conductual, que antes dependía de que el consumidor la declarara y hoy se registra sola.
\item Inventó variables de segmentación que no existían antes, como la frecuencia, el monto y el momento de las compras.
\end{enumerate}
\vspace{2pt}
\noindent\textbf{21.} La Figura 1 muestra la tasa de adopción de tres canales digitales simulados con el modelo de Bass (A: $p=0.08$, $q=0.10$; B: $p=0.03$, $q=0.38$; C: $p=0.01$, $q=0.55$). ¿Qué canal tiene el pico más tardío y más alto, y cómo se interpreta?
\vspace{-3pt}
\begin{enumerate}[label=\alph*),leftmargin=1.1em,itemsep=0pt,topsep=1pt,parsep=0pt,partopsep=0pt]
\item El canal C: con imitación dominante la adopción arranca lenta y despega después, con un pico más alto por el contagio social.
\item El canal B: es el único con balance entre innovación e imitación, y por eso su pico es el más tardío y más alto.
\item El canal A: como su coeficiente $p$ es el mayor, su pico llega al final del periodo y es el más alto de los tres.
\end{enumerate}
\vspace{2pt}
\noindent\textbf{22.} Un equipo estima que el patrón de adopción de su canal se parece a la curva C de la Figura 1. ¿Qué decisión de lanzamiento es más coherente?
\vspace{-3pt}
\begin{enumerate}[label=\alph*),leftmargin=1.1em,itemsep=0pt,topsep=1pt,parsep=0pt,partopsep=0pt]
\item Impulsar el contagio social (referidos, prueba social), porque el crecimiento depende sobre todo de la imitación.
\item Concentrar todo el presupuesto en publicidad masiva el primer mes, porque la adopción es casi toda independiente.
\item Esperar a que el pico ocurra sin intervenir, porque ni $p$ ni $q$ pueden influirse con acciones de mercadotecnia.
\end{enumerate}
\vspace{2pt}
\noindent\textbf{23.} Marcela compara precios, lee reseñas y revisa especificaciones antes de comprar, pero cambia de teléfono cada año y desecha el anterior. Según el capítulo, su consumo es:
\vspace{-3pt}
\begin{enumerate}[label=\alph*),leftmargin=1.1em,itemsep=0pt,topsep=1pt,parsep=0pt,partopsep=0pt]
\item Sostenible, porque comparar precios y leer reseñas antes de comprar reduce por sí solo el impacto ambiental de cada compra.
\item Inteligente, porque decide informada por datos, pero no necesariamente sostenible: el impacto ambiental es un criterio distinto.
\item Ni inteligente ni sostenible, porque cambiar de teléfono cada año invalida por completo la comparación que hizo antes de comprarlo.
\end{enumerate}
\vspace{2pt}
\noindent\textbf{24.} Una tienda de ropa pide a sus clientes que llenen, de forma voluntaria, un cuestionario de talla y estilo para recomendarles prendas. ¿Qué tipo de datos obtiene?
\vspace{-3pt}
\begin{enumerate}[label=\alph*),leftmargin=1.1em,itemsep=0pt,topsep=1pt,parsep=0pt,partopsep=0pt]
\item Datos de cero parte: el consumidor los entrega de forma voluntaria y explícita.
\item Datos de primera parte observados: la tienda los deduce de los clics y compras del cliente sin que él los declare.
\item Datos de tercera parte: los comparte un tercero sobre el comportamiento del cliente en otros sitios.
\end{enumerate}
\vspace{2pt}
\noindent\textbf{25.} Una app de entregas descubre, con analítica de embudo, que la mayoría de quienes abandonan lo hacen en la pantalla de pago. Corregir ese paso ilustra principalmente la ventaja de:
\vspace{-3pt}
\begin{enumerate}[label=\alph*),leftmargin=1.1em,itemsep=0pt,topsep=1pt,parsep=0pt,partopsep=0pt]
\item Reducir la fricción en el recorrido: detectar en qué punto se abandona el proceso de decisión y corregirlo.
\item Personalizar la oferta: la app ajusta precio y mensaje al perfil inferido de cada usuario para aumentar la conversión.
\item Anticipar tendencias: la app se adelanta a un cambio en los gustos y las preferencias del consumidor antes que la competencia.
\end{enumerate}
\vspace{2pt}
\noindent\textbf{26.} Una tienda opera sitio web, app y tiendas físicas, pero cada canal guarda a sus clientes en una base distinta y nadie reconoce que son la misma persona. Según Verhoef, Kannan e Inman (2015), esa tienda es:
\vspace{-3pt}
\begin{enumerate}[label=\alph*),leftmargin=1.1em,itemsep=0pt,topsep=1pt,parsep=0pt,partopsep=0pt]
\item Omnicanal, porque la tienda está presente en múltiples canales y puntos de contacto con el mismo cliente.
\item Omnicanal, siempre que cada canal optimice por separado su propio desempeño y su propia experiencia.
\item Multicanal, no omnicanal: falta la gestión sinérgica con una identidad única de cliente que atraviese los canales.
\end{enumerate}
\vspace{2pt}
\noindent\textbf{27.} Un directivo afirma: «ya no hay que preocuparse por las cookies de terceros, porque desaparecieron de los navegadores». ¿Cómo se valora esa afirmación a septiembre de 2026, según el capítulo?
\vspace{-3pt}
\begin{enumerate}[label=\alph*),leftmargin=1.1em,itemsep=0pt,topsep=1pt,parsep=0pt,partopsep=0pt]
\item Es correcta: Google las eliminó en 2024 y las sustituyó con las APIs de Privacy Sandbox, que hoy las reemplazan por completo.
\item Es incorrecta, y por eso la recolección de datos de primera y cero parte dejó de ser relevante para la estrategia.
\item No es verificable: Google revirtió su plan en 2024 y las cookies de terceros siguen activas por defecto en Chrome.
\end{enumerate}
\vspace{2pt}
\noindent\textbf{28.} Una empresa decide cuánto invertir en retener a un cliente comparándolo con lo que se espera que el cliente le deje. Esa cifra esperada corresponde al:
\vspace{-3pt}
\begin{enumerate}[label=\alph*),leftmargin=1.1em,itemsep=0pt,topsep=1pt,parsep=0pt,partopsep=0pt]
\item CLV (valor de vida del cliente): el ingreso esperado que el cliente genera durante toda su relación con la organización.
\item Ticket promedio: lo que el cliente gasta en una sola compra, sin considerar cuántas veces vuelve a comprar.
\item Costo de adquisición: lo que costó atraer a ese cliente la primera vez, sin importar lo que compre después.
\end{enumerate}
\vspace{2pt}
\noindent\textbf{29.} Un proveedor de analítica vende su plataforma con la promesa central de «anticipar tendencias antes que la competencia». ¿Qué observación hace un analista crítico, con base en el nivel de evidencia del capítulo?
\vspace{-3pt}
\begin{enumerate}[label=\alph*),leftmargin=1.1em,itemsep=0pt,topsep=1pt,parsep=0pt,partopsep=0pt]
\item Que es imposible de lograr, porque los datos de comportamiento no permiten ninguna clase de predicción útil sobre el consumidor.
\item Que es la ventaja con la evidencia más sólida del cuadro, pues la literatura arbitrada la respalda desde los trabajos de Smith (1956).
\item Que es la ventaja menos sostenida por evidencia: un patrón histórico no garantiza que se repita y hay riesgo de sobreajuste a ruido.
\end{enumerate}
\vspace{2pt}
\noindent\textbf{30.} ¿Cuál afirmación describe mejor la relación entre Smith (1956) y el proceso STP (segmentación, \textit{targeting}, posicionamiento)?
\vspace{-3pt}
\begin{enumerate}[label=\alph*),leftmargin=1.1em,itemsep=0pt,topsep=1pt,parsep=0pt,partopsep=0pt]
\item Smith propuso el STP completo en 1956 y los manuales posteriores solo lo difundieron entre profesionales sin modificarlo.
\item Smith aporta la base conceptual; el STP es su sistematización posterior en los manuales de Kotler y sigue vigente como marco.
\item El STP quedó obsoleto con la analítica digital, que lo reemplazó por la segmentación difusa como marco de decisión.
\end{enumerate}
\vspace{2pt}
\end{multicols}
\end{document}
